\documentclass [11pt,a4paper,oneside,openany]{book} %Classe del documento, formato carta, singola facciata, apertura capitoli pag destra/sinista (ininfluente se impostato oneside). Formato libro
\usepackage[italian]{babel} %Lingua Documento (da impostare per la sillabazione)
\usepackage{graphicx} %Pacchetto necessario per la gestione delle immagini
\usepackage{xcolor,colortbl} %Colori
\usepackage{enumitem} % pacchetto per gestire le enumerazioni e resume
\usepackage{fancyhdr} %Pacchetto per la gestione accurata della pagina
\usepackage{setspace} %Pacchetto necessario per i comandi successivi onehalfspacing, singlespacing....
\usepackage{titlesec} % Pacchetto per le sezioni e sottosezioni
\usepackage{longtable} %Pacchetto per la gestione delle tabelle grandi
\usepackage[colorlinks=true]{hyperref} %Segnalibri nel pdf finale, sotto relativi parametri
\usepackage[T1]{fontenc}
\usepackage[utf8]{inputenc} %Pacchetti per le lettere accentate utf8
\usepackage{listings} %Pacchetto per formattare codice sorgente
\usepackage{mathtools} %Pacchetto matematica
\usepackage{float} %Pacchetto per il posizionamento delle figure
\usepackage{amsthm}
\usepackage{supertabular}

\hypersetup{
	bookmarksnumbered=true,
	linkcolor=black,
	citecolor=black,
	urlcolor=black,
}

\usepackage{geometry} % Dimensione pagina

\geometry{a4paper} % Formato carta
\addtolength{\textheight}{75pt} %Margini
\oddsidemargin 30pt

\setcounter{secnumdepth}{2}
\setcounter{tocdepth}{3}

\newcommand{\RASDref}[1]{\textbf{RASD: #1}}
\newcommand{\DDref}[1]{\textbf{DD: #1}}
\newcommand{\classDescription}[1]{\subsubsection{#1}\label{class:#1}}
\newcommand{\methodDescription}[1]{\textit{#1}(..)}


\newcommand{\testTable}[3]{
\begin{longtable}{|p{14cm}|} \hline
\textbf{Precondition and action tested: } #1  \\ \hline
\textbf{Expected behaviour: } #2  \\ \hline
\textbf{Actual outcome: } #3  \\ \hline
\end{longtable}
}

% #1: name of the image on the images folder. eg: middleware.png
\newcommand{\testImage}[1]{
\begin{figure}[H]
\begin{center}
\includegraphics[width=250pt, keepaspectratio]{"images/"#1}
\caption{Graphical representation of the actual outcome}
\end{center}
\end{figure}
}

\newcommand{\TestTable}[4]{
\begin{longtable}{|p{14cm}|} \hline
\textbf{Precondition: } #1  \\ \hline
\textbf{Action tested: } #2  \\ \hline
\textbf{Expected behaviour: } #3  \\ \hline
\textbf{Actual outcome: } #4  \\ \hline
\end{longtable}
}

\newcommand{\TestImage}[2]{
\begin{figure}[H]
\begin{center}
\includegraphics[width=250pt, keepaspectratio]{"images/#1".png}
\caption{#2}
\end{center}
\end{figure}
}

\definecolor{darkgreen}{rgb}{0,0.6,0}
\lstset{frame=tb,
  language=php,
  aboveskip=3mm,
  belowskip=3mm,
  showstringspaces=false,
  columns=flexible,
  basicstyle={\small\ttfamily},
  numbers=none,
  keywordstyle=\color{blue},
  commentstyle=\color{black},
  breaklines=true,
  breakatwhitespace=true,
  tabsize=3
}

\begin{document} %Inizio Documento

%%FRONTESPIZIO%%=========================================
\begin{titlepage}

 \begin{center} 
     \includegraphics[height=6cm]{logo.eps}\\
     \vspace{4em}
     {\Large \textsc{Software Engineering II}}\\
     \vspace{6em}
     {\LARGE \textbf{Travlendar+}}\\
     \vspace{3em}
     {\Large \textsc{Implementation \& Testing}}\\
     \vspace{1em}
     {\Large \textsc{Document}}\\
 \end{center}
 
    \vskip 1cm
 
 	Authors:
 	\vspace{0.5em}
 	\begin{center}
      {\Large \textit{Edoardo D'Amico}}\\
      {\Large \textit{Gabbolini Giovanni}}\\
      {\Large \textit{Parroni Federico}}\\
    \end{center}

	\vskip 1.5cm    
    
    \begin{center}
    Git Repository: \url{https://github.com/keyblade95/DamicoGabboliniParroni}
	\end{center}

\vskip 2cm
\begin{center}
{\normalsize 1$^{st}$ December 2017}
\end{center}

\end{titlepage}

\newpage

%%INTESTAZIONI PAGINA%%==================================
\pagestyle{fancy}
\renewcommand{\chaptername}{Chapter}
\renewcommand{\chaptermark} [1]{\chaptername\ \thechapter.\ #1}{} 
\renewcommand{\chaptermark}[1]{\markboth{\thechapter.\ #1}{}} 
\renewcommand{\sectionmark}[1]{\markright{\thesection\ #1}}
\fancyhf{}
\fancyhead[LE,RO]{\bfseries\thepage} 
\fancyhead[LO,RE]{\bfseries\leftmark}
\fancypagestyle{plain}{%
\fancyhead{} % leva l'intestazione
\renewcommand{\headrulewidth}{0pt} % e la linea
}

%%INDICE%%==============================================
\renewcommand{\contentsname}{Indice}
\renewcommand{\listfigurename}{Elenco delle Figure}
\renewcommand{\listtablename}{Elenco delle Tabelle}
\tableofcontents

\linespread{1.5}\selectfont		%interlinea
%%CAPITOLI===============================================
%%CAPITOLO 1: Introduzione=======================================
\chapter{Introduction}

\section{Purpose}
The purpose of this document is to provide more technical details than the RASD about Travelendar+ system. This document is addressed to developers and aims to identify:
\begin{itemize}
\item The high level architecture
\item The design patterns
\item The main components and their interfaces provided one for another
\item The Runtime behavior
\end{itemize}

\section{Scope}

Travlendar+ is an application that aims to ease the management of the daily appointments of a registered user by providing a desired schedule for those. 
The application runs within the context of a restricted geographic area and thus should consider only the travel means available in this scope.
The application is able to interact with different data sources, in particular Travel Means APIs and User device GPS, and to arrange the appointments by exploiting these information.
Once the schedule is computed, the application gives to the user the possibility to buy tickets for the involved travel means, then he/she can run it and follow directions provided by the application. 

\section{Definitions, Acronyms, Synonims}

\subsection{Definitions}
\theoremstyle{definition}
\newtheorem{definition}{Definition}[section]
 
\begin{definition} \label{def:userdata}
The user data include:
\begin{itemize}
\item Schedules;
\item Appointments;
\item User parameters.
\end{itemize}

\end{definition}
  
\begin{definition} \label{def:device}
A device is a PC, a Tablet or a Smartphone in which run the last version of his O.S.;
\end{definition}

\begin{definition} \label{def:schedule}
A Schedule is a set of time-ordered and not overlapping appointments where their starting times are fixed and they're linked to each other by a path covered with a specific transportation mean;
\end{definition}

\begin{definition} \label{def:bestTravelMean}
The bestTravelMean is the travel mean that satisfies the constraints on the appointment and on the schedule and that optimize the chosen criteria.
\end{definition}


\subsection{Acronyms}
%Acronyms used in the text:
\begin{itemize}
\item API: Application Programming Interface;
\item CRUD: Create/Read/Update/Delete;
\end{itemize}

%\reqref{req:R1}
%\goalref{goal:G1}

\section{Revision history}

\section{Reference Documents}
RASD document

\section{Document structure}
\begin{itemize}
\item Introduction: this section introduces the design document. It contains a justification of his utility
\item Architecture Design: this section is divided into different parts:
\begin{enumerate}
\item Overview: this section explains the division in tiers of our application and lists further details that were not specified in the RASD.
\item Components View: this sections gives a global view of the components of the application, other than a brief description, and describes how they communicate
\item Deployment view: this section shows the components that must be deployed to have the application running correctly
\item Runtime view: sequence diagrams are represented in this section to show how the components interact and collaborate during the fullfilling of a specific task. Each task is associated to each goal specified on the RASD document
\item Component interfaces: the interfaces between the components are presented in this section
\item Selected architectural styles and patterns : this section explains the architectural and design choices taken during the creation of the application
\end{enumerate}
\item Algorithms Design: this section describes the most important algorithmic parts. Pseudo code is used in order to hide unnecessary implementation details and abstract from a specific programming language
\item Requirements Traceability: this section aims to explain how the decisions taken in the RASD are linked to design elements
\end{itemize}



 

\chapter{Project Analized}
\chapter{Installation Instruction}
The software has been installed trough the apk whitin the deliverables folder in the github repository how said from the group in the ITD document. The application has been installed on two different virtual devices, the nexus 5 emulator, with android version 8.0(Oreo) and in the nexsus 5x with the same version of android, that has been done for test how the application can adapt itsel to different display size, 5 for the first one and 5.5 for the second one. There were no inchorency on the installation instruction released on the ITD apart for the presence of the application in the google playstore, in fact the application can't be downloaded from it.
\chapter{Tests}

1. non viene aggiornata la view dopo l'inserimento di un appointment\\
2. i colori degli appuntamenti cambiano\\
3. la fattibilità di uno schedule non è controllata\\
4. non è controllata la consistenza delle location\\
5. se viene cambiato numero di bike sharing o owned bike crasha\\
6. se l'ora di inizio < ora fine crasha\\
7. se si mettono appuntamenti sovrapposti l'applicazione crasha\\
8. in un display da 5' il layout non permette di interagire con alcuni elementi grafici\\

\section{Test of the implemented functionalities}
The tests has been done on the implemented functionalities declared on the ITD document released by the other team, for each of them at least on test has been carried out.

\subsection{Insertion of events and free times}

The insertion of an event can be done by clicking the button the main page of the application

\TestImage{func1/MainView}{Main page view}

the user is redirect to an event creation view with some fields to fill as shown below.

\TestImage{func1/AddEventPage}{Event Creation View}

The view has some missing graphic constraints that put the two clock picker in the middle of the view hiding the bar that made possible the selection of a travel mean for the event that is currently being created.\\
Another small issue is that the year of the selected day is wrong(is 0118 instead of 2018).\\
There isn't any check for the feasibility of the location inserted.\\
\subsubsection{Test1}
A first test has been carried out with a normal insertion of an event without free time and without repetition.
After the click of the add event button the user is redirected to the main page but the appointment has not been added,although with a change of date on the calendar (change date and after rechange on the current date) the appointment is added, as shown in the image below,probably a little problem in refreshing the view after the insertion.

\TestImage{func1/AppAdded1}{insertion of an event}

\subsubsection{Test2}
here is tryed to add an event with repetition that is possible pushing the button on the event creation view named "repetition" the user is redirected to the following view

\TestImage{func1/RepetitionApp}{Repetition view}

the confirm button has some missign constraints so it stay on the middle of the field for the insertion of the date.\\
After the click of the confirm button the user is redirected to the main page after the refresh of the view the appointment is visible only on the first date of the slot selected for the repetition, on the other days in which the appointment should be shown nothing has been added.

\subsubsection{Test3}
Here is tryed to add two overlapping appointments, in this case after the insertion of the second one the application crash.

\subsubsection{Test4}
Here is tryed to add a free time, when the save button is clicked the application crash.

\subsection{Notification System}
This functionality is declared to be implemented on the ITD document but it isn't working

\subsubsection{SettingPreferences}

\subsubsection{Persistance Database to store events and free times}

\subsubsection{BikeMI API}
This functionality is declared to be implemented on the ITD document but it isn't working

\subsubsection{Wheather API}
This functionality is declared to be implemented on the ITD document but it isn't working

\subsubsection{preferences management}





\chapter{Additional Points}
\chapter{Effort Spent}
\begin{itemize}
\item Federico Parroni: \textbf{130 hours};
\item Edoardo D'Amico: \textbf{130 hours};
\item Giovanni Gabbolini: \textbf{130 hours}.
\end{itemize}

%%APPENDICI==============================================
\appendix
%%\input{appendiceA}

\end{document}